% !TEX root = ../../main.tex

\section{Introduction}

Hierarchical clustering, meaning a family of methods that builds nested groups
by repeatedly merging similar samples or sample groups, is attractive when the
analyst wants more than a flat partition. In this setting, a tree is a rooted
diagram of nested groups of
samples. The leaves of the tree are the observed samples themselves. As the
clustering algorithm merges similar samples or groups, it creates internal
nodes, each of which represents the union of all leaves below it. The root is
the final node containing the full data set. Such a representation exposes
coarse-to-fine structure, keeps near-miss cluster boundaries visible, and makes
it possible to discuss not only the final grouping but also the intermediate
grouping decisions that led to it. In practice, however, most hierarchical
workflows still cut the tree using geometric rules such as a fixed
height, an inconsistency threshold, or a manually chosen number of clusters.
Those choices are often difficult to defend when the goal is not only
descriptive grouping but also reproducible statistical claims.

Tree-Break Selection takes a different view of this tree. An internal node is not merely a
merge produced by the algorithm; it is a candidate decision point. If a parent
node $u$ has children $l(u)$ and $r(u)$, then that parent proposes a possible
split of the data into two child subtrees. The central question is whether this
proposed split should survive in the final partition. A linkage tree alone
cannot answer that question. Here, a linkage tree means the hierarchy returned
by the agglomerative merging rule itself. It proposes candidate splits, but it
does not tell
us whether at least one child genuinely departs from the parent's feature
distribution, whether the two children are actually distinct enough to
represent separate clusters, or whether scanning the entire tree has produced
false positives.

Tree-Break Selection treats that split decision as a two-stage statistical question. A split
must first show evidence that at least one child departs from its parent. It
must then show that the two children are distinct from one another. The final
partition is built only after these local decisions have been corrected for the
many candidate splits examined across the hierarchy.

Four inferential complications shape the design. The child is a subset of its
parent, so the edge comparison is nested rather than an ordinary two-sample
test. The edge evidence is high-dimensional and correlated, so a direct sum of
squared feature scores would overcount repeated or noisy directions. Many candidate edges
are examined across the tree, so valid local $p$-values still require control
for many tests. Finally, sibling comparisons are evaluated on a hierarchy
constructed from the same data, so the raw sibling statistics are selected and
inflated. Tree-Break Selection is one response to these four problems.

This goal is related to, but distinct from, existing approaches for attaching
uncertainty statements to hierarchical clustering. Multiscale-bootstrap methods
such as \texttt{pvclust} estimate cluster-level uncertainty by repeatedly
rebuilding dendrograms under resampling \citep{suzuki2006pvclust}. Sequential
hierarchical clustering significance procedures test splits while controlling
tree-aware error rates under explicit null models \citep{kimes2017statistical}.
Recent selective inference work conditions on the clustering event and computes
valid post-selection p-values for selected hierarchical clusters under Gaussian
mean models \citep{gao2024selective}.

Tree-Break Selection does not currently derive an exact selective law for its tree-selected
sibling tests. Instead, it implements a projected-Wald testing pipeline over an
explicit feature-space contract, with an empirical-null inflation layer for the
sibling stage. The restored rule uses the estimated scale in a plug-in chi-square tail
when internal support exists; selected-tail validity remains unproven. That difference is
important: the present draft should be read as a methods specification and
validation plan, not as a proof of exact selective Type~I error control.

Tree-Break Selection works on typed feature blocks rather than one universal binary
representation. Bernoulli coordinates enter directly; categorical variables
enter as one-hot blocks but are tested with their drop-last multinomial
covariance. Explicit continuous inputs enter as empirical-Gaussian blocks whose
subtree covariances are estimated from descendant leaves. Discretized Gaussian
representations remain benchmark variants rather than the continuous-data
contract.
Once the representation is fixed, the inferential objects are no longer
individual samples but subtrees, so each node carries a feature-wise
distribution estimated from its descendant leaves. Statistical test annotations
are computed before the actual tree walk that creates cluster assignments. This
separation between evidence calculation and final traversal is useful both
computationally and scientifically: the statistical support for each potential
split is stored explicitly and can be audited independently of the resulting
partition.

The presentation follows that decision problem in order. It first defines the
tree and node distributions. It then gives the child-parent edge test, the
tree-level multiplicity correction, the sibling test with its post-selection
inflation layer, and the final top-down traversal that converts test
annotations into clusters.
